% Study sheet: Volume II, Chapter 15, Sustainable AI.
% Numbers: \val{...} from sustainable_ai_studysheet_numbers.tex, written by sustainable_ai_studysheet_numbers.py.
\documentclass{studysheet}
\input{sustainable_ai_studysheet_numbers}
\sheetvolume{2}
\sheetchapter{15}{Sustainable AI}
\sheetsources{Numbers from mlsysim and the chapter's scenarios and data; energy coefficients are the chapter's modeled values; assumed inputs marked.}
\newcommand{\cotwo}{CO\textsubscript{2}}

\begin{document}
% ===================================================================== FRONT: LEARN
\frontside
\begin{bigidea}
A fleet computes only as fast as its site delivers power and removes heat, and every joule carries
carbon set by where and when it runs, plus the carbon of building the hardware. Measure under a declared
\hl{boundary}, find the \hl{largest term}, pull the lever that changes it. In C\textsuperscript{3} terms,
compute, communication and coordination share one physical budget.
\end{bigidea}

\begin{leftcol}
\sheetsec{Key ideas}
\begin{ideas}
\idea{Power is a ceiling; energy is an account.}{Draw may never pass the feed; energy integrates power
  over time, so a short burst can be infeasible at modest energy.}
\idea{A footprint needs its contract.}{Functional unit, boundary, attribution rule, carbon signal, time
  and place: what one study counted moved a prompt's energy \val{qbound}$\times$.}
\idea{Operational carbon is a chain.}{IT energy times PUE times grid factor. PUE counts only
  overhead; evaporative cooling lowers it by spending water.}
\idea{Carbon starts before boot.}{Manufacturing is spent before power-on; its share grows on clean
  grids and at low use.}
\idea{Moving work changes the factor, not the energy.}{Use a marginal forecast; a pause must cut metered
  power and beat checkpoint, idle and restart carbon by the deadline.}
\idea{Demand outruns efficiency.}{Fleet energy is energy per job times jobs; scaling and rebound grow
  jobs, so set absolute budgets.}
\end{ideas}

\vspace{10pt}
\sheetsec{Where each lever acts}
{\centering
\begin{tikzpicture}[sheet]
  \node[accbox, minimum width=58pt, minimum height=20pt] (it) at (30,52) {IT energy};
  \node[accbox, minimum width=58pt, minimum height=20pt] (fac) at (125,52) {facility energy};
  \node[accbox, minimum width=58pt, minimum height=20pt] (op) at (220,52) {operational \cotwo};
  \node[nbox, minimum width=58pt, minimum height=20pt] (mfg) at (30,0) {manufacturing};
  \node[hotbox, minimum width=58pt, minimum height=20pt] (life) at (220,0) {lifecycle \cotwo};
  \draw[arr] (it) -- node[lab, above=3pt] {$\times$ PUE} node[note, below=4pt] {cooling} (fac);
  \draw[arr] (fac) -- node[lab, above=3pt] {$\times$ grid} node[note, below=4pt] {place, hour} (op);
  \draw[arr] (mfg) -- node[lab, above=3pt] {per device-day} node[note, below=4pt] {longer life} (life);
  \draw[arr] (op) -- node[lab, right=3pt] {$+$} (life);
\end{tikzpicture}\par}

\vspace{14pt}
\sheetsec{Numbers to know}
\begin{sheettable}{@{}l L@{}}
US-average grid & \val{usci} g/kWh, facility PUE \val{uspue}\\
Quebec hydro grid & \val{qcci} g/kWh\\
Embodied, per device & A100 \val{aembodied} kg, H100 \val{hembodied} kg; life \val{lifestd} years (\val{lifeext} extended)\\
Modeled GPU energy & \val{eop} pJ/FLOP; DRAM \val{ebyte} pJ/byte\\
DGX H100 rack & GPUs \val{rackgpu} kW; whole IT load \val{rackitload} kW\\
\end{sheettable}
\end{leftcol}%
\begin{rightcol}
\sheetsec{Formulas you will use}
\fcard{Power ceiling and run energy}%
  {\FacilityPower = \text{PUE}\times\InformationTechnologyPower \le \GridFeedCapacity
   \qquad \TotalEnergy = \AveragePower\times\JobDuration}%
  {Feed $\GridFeedCapacity$ caps facility draw $\FacilityPower$; IT power $\InformationTechnologyPower$
   includes hosts, network and conversion, not only TDP. Energy is average power $\AveragePower$ over
   run time $\JobDuration$.}%
  {Step energy is fleet power times the Fleet Law step time: exposed communication and waits cost joules.}{}
\fcard{Carbon account}%
  {\begin{gathered}\text{CO}_{2,\text{op}} = \under{\InformationTechnologyEnergy\times\text{PUE}}{facility energy}
   \times\text{CO}_{2,\text{grid}}\\[2pt]
   \text{CO}_{2,\text{emb}} = \frac{\text{CO}_{2,\text{mfg}}}{\val{daysperyear}\,\HardwareServiceLife}\times\text{device-days}
   \end{gathered}}%
  {The grid factor (g/kWh) converts facility electricity to carbon; manufacturing is charged per
   device-day over a life $\HardwareServiceLife$ in years.}%
  {Lifecycle = training (times development runs) + inference (recurs) + embodied.}{}
\fcard{Energy crossover}%
  {\ArithmeticIntensity = \frac{\Operations}{\DataVolume}
   \qquad \EnergyCrossoverIntensity = \frac{\EnergyPerByteMoved}{\EnergyPerOperation}}%
  {Intensity $\ArithmeticIntensity$: FLOPs $\Operations$ per byte $\DataVolume$. Below
   $\EnergyCrossoverIntensity$, moving bytes ($\EnergyPerByteMoved$ each) costs more than arithmetic
   ($\EnergyPerOperation$ per FLOP). In FP16 decode one two-byte weight load serves $\BatchSize$
   tokens of two FLOPs each, so $\ArithmeticIntensity = \BatchSize$.}%
  {The performance ridge sits higher: between the two, energy is arithmetic-heavy but time memory-paced.}{}
\end{rightcol}

% ===================================================================== BACK: PRACTICE
\backside
\begin{answers}
\ans{Finish}{\val{fopqc} + \val{wemb} = \val{ftotal} kg; embodied \val{wemb} $\div$ \val{ftotal} = \val{fsharepct}\%.
  Operational fell \val{fratio}$\times$, manufacturing not at all: electricity still leads, but
  service life and utilization now matter (``before boot'').}
\ans{1}{(a) \val{qprompts} million $\times$ \val{qfull} Wh = \val{qdayfull} kWh a day; \val{wfac} $\div$
  \val{qdayfull} = \val{qdaysfull} days. (b) \val{qdaynarrow} kWh a day: \val{qdaysnarrow} days,
  \val{qbound}$\times$ later: the narrow figure drops hosts, idle machines, overhead and all but
  the best sites (contract idea).}
\ans{2}{(a) $\ArithmeticIntensity$ = \val{dint}: \val{bpp} $\times$ \val{ebyte} = \val{dmove} pJ moved against
  \val{fpp} $\times$ \val{eop} = \val{dcomp} pJ, \val{dshare}\% movement. (b) One load serves \val{qbatch} tokens:
  \val{qmove} against \val{dcomp} pJ, \val{qshare}\%. (c) Batch \val{balance}; time stays memory-paced below
  the higher performance ridge (crossover card).}
\columnbreak
\ans{3}{Saves \val{rtdp} W $\times$ \val{rsave}\% = \val{rsavedkw} kW, \val{rkwh} kWh a day. US: \val{rkgus} kg a
  day, \val{hembodied} $\div$ \val{rkgus} = \val{rdaysus} days: replace. Quebec: \val{rkgqc} kg a day,
  \val{rdaysqc} days (\val{ryearsqc} years), past a \val{lifestd}-year life: keep (``before boot'').}
\ans{4}{(a) A marginal hourly forecast; an annual average has no hours. (b) Checkpoint, idle held cluster, migration and restart, within the deadline.
  (c) Hourly facility kWh times marginal factor over the whole job, shifted against unshifted (``moving work'').}
\end{answers}

\begin{leftcol}
\sheetsec{Worked example}
\begin{practicebox}\boxtitle{What does a \val{wdays}-day run on \val{wgpus} A100s emit, and how much is manufacturing?}
\val{wgpus} A100s for \val{whours}~h, \assumed{} at their rated \val{wtdp}~W throughout (an upper
bound); the DGX H100 rack's IT ratio \val{itratio} (hosts, network, conversion) as a proxy;
PUE \val{uspue}; US-average grid \val{usci}~g/kWh; \val{aembodied} kg per A100 over a
\val{lifeext}-year extended life.\par\vspace{2pt}
\begin{steps}
\item Accelerators: \val{wgpus} $\times$ \val{wtdp}~W $\times$ \val{whours}~h = \val{wacc}~kWh.
\item IT: $\times$ \val{itratio} = \val{wituse}~kWh; facility: $\times$ \val{uspue} = \res{\val{wfac}~kWh}.
\item Operational: \val{wfac} $\times$ \val{usciskg}~kg/kWh = \res{\val{wopus}~kg} (\val{woptons}~t).
\item Embodied: \val{aembodied} $\div$ (\val{lifeext} $\times$ \val{daysperyear}) = \val{wperday} kg per
  device-day; $\times$ \val{wdevdays} device-days = \val{wemb} kg, \res{\val{wsharepct}\%} of the total.
\end{steps}
\takeaway{On this grid, electricity dominates: cut energy or the factor.}
\end{practicebox}

\vspace{10pt}
\sheetsec{Your turn: finish it}
\begin{practicebox}\boxtitle{Same run, hydro grid: does manufacturing still not matter?}
\textbf{Given:} the same facility energy on Quebec's \val{qcci}~g/kWh grid: \val{wfac} $\times$
\val{qccikg} = \val{fopqc} kg operational.\par
\textbf{You:} find the total, the embodied share and the lever that matters now.\par\vspace{72pt}
\end{practicebox}
\end{leftcol}%
\begin{rightcol}
\sheetsec{Common mistakes}
\mistake{``Budget the site by GPU TDP.''}%
  {\assumed{} a \val{feed} MW feed at PUE \val{uspue} allows \val{itcap} MW of IT: \val{rackstdp} racks
   by GPU power (\val{rackgpu} kW), only \val{racksit} at the whole IT load (\val{rackitload} kW).}{}
\mistake{``The lower-PUE site is greener.''}%
  {\assumed{} per IT kWh, PUE \val{pueliq} on a \val{iowaci} g coal-gas grid emits \val{pga} g; PUE
   \val{pueleg} on a \val{franceci} g nuclear grid, \val{pgb} g: \val{pgratio}$\times$ less.}{}
\mistake{``Cutting energy per query \val{unitgain}$\times$ cuts our footprint.''}%
  {\assumed{} price and energy per query fell \val{unitgain}$\times$, volume grew \val{volgrow}$\times$: total energy
   rose \val{netgrow}$\times$.}{}

\vspace{4pt}
\sheetsec{Check yourself}
\begin{qs}
\item \assumed{} \val{qprompts} million prompts a day at the measured \val{qfull} Wh full stack (\val{qnarrow} Wh:
  active accelerators in the most efficient sites). (a) After how many days does serving equal the worked example's
  \val{wfac} kWh? (b) With the accelerator-only figure? Why the gap?
\item FP16 decode. (a) At batch one, what share of energy per token is
  movement? (b) At batch \val{qbatch} (assumed)? (c) Where do the terms balance? Why can time stay
  memory-bound?
\item \assumed{} accelerator power only, \val{lifestd}-year life: a replacement with an H100's manufacturing carbon needs \val{rsave}\% less energy than an
  H100 running continuously at \val{rtdp} W. When does it repay manufacturing on the US-average grid,
  and on Quebec's?
\item A scheduler pauses a checkpointed run each evening, choosing hours by the annual-average
  factor. (a) Which signal? (b) What must the avoided carbon exceed? (c) What
  measurement proves a saving?
\end{qs}
\end{rightcol}
\end{document}
